\documentclass[conference]{IEEEtran}
\IEEEoverridecommandlockouts
\usepackage{cite}
\usepackage{amsmath,amssymb,amsfonts}
\usepackage{algorithmic}
\usepackage{graphicx}
\usepackage{textcomp}
\usepackage{xcolor}
\usepackage{url}
\def\BibTeX{{\rm B\kern-.05em{\sc i\kern-.025em b}\kern-.08em
    T\kern-.1667em\lower.7ex\hbox{E}\kern-.125emX}}

\begin{document}

\title{Machine Unlearning in the Era of Large Language Models: Challenges and Methodologies}

\author{\IEEEauthorblockN{Kevin Kayasit Lendvay}
\IEEEauthorblockA{\textit{Department of Information Systems} \\
\textit{Eötvös Loránd University}\\
Budapest, Hungary \\
ajz5ct@inf.elte.hu}
}

\maketitle

\section{Introduction}
Machine learning models, particularly Large Language Models (LLMs), have recently seen an unprecedented proliferation and are 
trained on massive datasets \cite{wang2026machineunlearningcomprehensivesurvey}. These datasets frequently encompass sensitive 
personal data, copyrighted materials, or even harmful and illegal information \cite{YAN2025100300, ICLR2025_4556f539}. Increasingly 
stringent data protection regulations, such as the "Right to be Forgotten" mandated by the GDPR, pose a novel and critical 
technological challenge for system operators: they must enable the effective removal of user data and its residual influence 
from already trained models \cite{wang2026machineunlearningcomprehensivesurvey, Liu2025}.

Traditionally, the complete removal of a specific data point's influence would necessitate the complete retraining of the 
model from scratch (i.e., naive retraining) \cite{wang2026machineunlearningcomprehensivesurvey}. However, in the context of 
contemporary models comprising billions of parameters, this procedure is computationally prohibitive, highly time-consuming, 
and often entirely infeasible \cite{7163042, Liu2025}. To address this challenge, the rapidly evolving research domain of machine 
unlearning has emerged, aiming to develop methodologies capable of removing the influence of targeted data without repeating the 
arduous training process \cite{wang2026machineunlearningcomprehensivesurvey, 7163042}. The motivation behind this field is 
twofold: mitigating privacy and security risks (e.g., memory poisoning or data leakage) and ensuring the overall safety and 
reliability of the models \cite{YAN2025100300}. The paramount importance of this area is evidenced by the fact that in 
recent years, machine unlearning algorithms, their verification, and their security limitations have become central 
research directions at top-tier artificial intelligence conferences (e.g., ICLR, ICML, S\&P).

\section{Related Work}
The literature on machine unlearning can be broadly classified into two primary categories based on the applied techniques 
and accuracy guarantees: exact unlearning and approximate unlearning methodologies \cite{wang2026machineunlearningcomprehensivesurvey, Liu2025}. 

Exact unlearning guarantees the complete statistical elimination of a given data point's influence from the model, 
rendering it as if the data had never been included in the training set \cite{wang2026machineunlearningcomprehensivesurvey}. 
A pioneering approach to this was introduced by Cao and Yang, who transformed learning algorithms into a statistical query (SQ) 
summation form \cite{7163042}. This transformation allows for the rapid deduction of individual data influences, achieving an 
asymptotically higher speed compared to traditional retraining \cite{7163042}. Another prevalent and effective exact 
method is the SISA (Sharded, Isolated, Sliced, and Aggregated) framework \cite{9519428}. SISA divides the dataset 
into isolated partitions (shards) and slices during training; thus, when deleting a data point, only a small constituent 
model requires retraining \cite{9519428}. Although SISA drastically accelerates unlearning, literature indicates that this 
speedup demands a trade-off: the procedure incurs a degradation in model accuracy and significant storage overhead \cite{9519428}.

For modern deep neural networks and LLMs, implementing exact unlearning is cumbersome due to high architectural complexity; 
consequently, approximate unlearning has become the predominant approach \cite{Liu2025}. Guo et al. developed a certified data 
removal mechanism for $L_2$-regularized linear models \cite{guo2023certifieddataremovalmachine}. Their method combines 
a Newton step with loss perturbation to achieve statistical indistinguishability akin to differential privacy, though they 
acknowledge that end-to-end unlearning in deep networks with non-convex loss functions remains an unresolved 
challenge \cite{guo2023certifieddataremovalmachine}. 

In the realm of large language models, approximate methods such as gradient ascent or Negative Preference Optimization (NPO) have 
introduced new challenges. Studies have demonstrated that these techniques often result in merely "spurious unlearning" due to the 
"squeezing effect". In such cases, while the probability of the target response is reduced, the model simply rephrases the harmful 
knowledge without actually erasing it \cite{ICLR2026_a63ce8e6}. The Bootstrapping (BS) framework by Li et al. was designed to 
address this phenomenon by suppressing the model's own high-probability generations at the token (BS-T) and sequence (BS-S) levels, 
thereby preventing the latent persistence of knowledge \cite{ICLR2026_a63ce8e6}. Furthermore, Liu et al. highlighted that 
LLM unlearning can be closely connected with model editing and adversarial training; hence, the focus of unlearning must 
extend beyond mere data deletion to encompass the removal of harmful model capabilities \cite{Liu2025}.

Currently, the evaluation of machine unlearning and the analysis of privacy threats targeting these algorithms have 
become critical research trajectories \cite{YAN2025100300, ICLR2025_4556f539}. Shi et al. introduced the 
MUSE (Machine Unlearning Six-Way Evaluation) framework, which assesses algorithms across six dimensions 
(e.g., verbatim and knowledge memorization, utility preservation, and scalability) \cite{ICLR2025_4556f539}. The MUSE benchmark 
revealed that although current algorithms (such as NPO or Task Vectors) can impede memorization, they drastically impair the 
model's general utility and are unsustainable in handling large-scale privacy leakage risks \cite{ICLR2025_4556f539}. 
This pessimistic yet realistic outlook is corroborated by recent comprehensive literature reviews. Yan et al. meticulously 
categorized active and passive privacy attacks targeting LLMs and autonomous LLM agents \cite{YAN2025100300}. 
Their research underscores that defense strategies against membership inference, model stealing, and memory poisoning attacks 
must be extended across the entire lifecycle of the models—from pre-training to inference \cite{YAN2025100300}.

\bibliographystyle{IEEEtran}
\bibliography{references}

\end{document}